\chapter{Averaging over polynomial moduli}\label{ch:averaging}

After \cref{ch:variance,ch:correlations} the variance is expressed through two sums over $k$: the density term $T_f(z,y)$ and the main part of the off-diagonal term,
\begin{equation}\label{eq:Mf}
M_f(x;z,y)=\sum_{z<k\le y}f'(k)\sum_{1\le h\le x/f(k)}\bigl(x-hf(k)\bigr)\Sing\bigl(hf(k)\bigr).
\end{equation}
Both involve arithmetic functions of $f(k)$. Expanding them over $r$ with the coefficients $\lambda(r)$ of \cref{ch:densities} makes the dependence on $k$ periodic modulo $r$. Partial summation then replaces the sum over $k$ by an integral in the variable $t=f(k)$, with $f'(k)$ acting as the Jacobian. This is the role of the weight $f'(k)$. (Archived source: proposition P4, a crux of the plan.)

\section{The averaged singular series}

For $h\ge1$ and a prime $p$ put
\[
\vartheta_p(h)=\frac1{p^2}\#\{a\bmod p^2:\ p^2\mid hf(a)\}.
\]
If $j=\max\{2-v_p(h),0\}$, then $p^2\mid hf(a)$ is equivalent to $p^j\mid f(a)$. Counting lifts from residues modulo $p^j$ to residues modulo $p^2$ gives
\begin{equation}\label{eq:deltap}
\vartheta_p(h)=\frac{\rho(p^{j})}{p^{j}}\in[0,1],\qquad\text{so}\qquad
1-\frac2{p^2}+\frac{\vartheta_p(h)}{p^2}=w_p(h),
\end{equation}
where $w_p(h)$ is the local factor in \eqref{eq:W-intro}. Comparing with the second form of \eqref{eq:sing}, $W_f(h)=\prod_pw_p(h)$ is what one obtains by replacing the indicator $\one_{p^2\mid hf(k)}$ in $\Sing(hf(k))$ by its average $\vartheta_p(h)$ over $k$. \Cref{lem:W-ident} below makes this precise.

\begin{lemma}\label{lem:W-basic}
For every $h\ge1$ and prime $p$, $1-2p^{-2}\le w_p(h)\le1-p^{-2}$. Hence the product $W_f(h)$ converges absolutely, uniformly in $h$, and $0<W_f(h)\le1$. The series
\[
\widetilde R_f=\sum_{r\ge1}\lambda(r)\frac{\rho(r)}r
\]
converges absolutely and $0\le\widetilde R_f\le c$.
\end{lemma}

\begin{proof}
The bounds for $w_p$ follow from $\vartheta_p(h)\in[0,1]$, so $|w_p(h)-1|\le2p^{-2}$, and \cref{lem:products} applies. The factors are positive, so $W_f(h)>0$, and they are at most one. Finally $0\le\rho(r)/r\le1$ and $\sum_r\lambda(r)=c$.
\end{proof}

For $r,h\ge1$ define
\begin{equation}\label{eq:beta}
\beta_r(h)=\frac1r\sum_{a=1}^r\frac{c_r(hf(a))}{\varphi(r)},
\end{equation}
the mean value of the $r$-periodic sequence
\[
k\longmapsto\frac{c_r(hf(k))}{\varphi(r)}.
\]
Note $|\beta_r(h)|\le1$ by \cref{lem:ramanujan}(b).

\begin{lemma}\label{lem:W-ident}
For every $h\ge1$, $\sum_{r\ge1}\lambda(r)\beta_r(h)=W_f(h)$, the series converging absolutely and uniformly in $h$.
\end{lemma}

\begin{proof}
Absolute and uniform convergence follow from $|\beta_r(h)|\le1$ and $\sum_r\lambda(r)=c$. If $(r,s)=1$, the Chinese remainder theorem identifies $a\bmod rs$ with pairs $(a\bmod r,\ a\bmod s)$. Moreover $c_{rs}=c_rc_s$, $c_r(hf(a))$ depends only on $a\bmod r$, and $\varphi(rs)=\varphi(r)\varphi(s)$. Hence $\beta_{rs}(h)=\beta_r(h)\beta_s(h)$. Thus $\lambda(r)\beta_r(h)$ is $c^2$ times a multiplicative function supported on cube-free $r$, and \cref{lem:euler} gives
\[
\sum_r\lambda(r)\beta_r(h)=c^2\prod_p\Bigl(1+\frac{(p-1)\beta_p(h)+(p^2-p)\beta_{p^2}(h)}{(p^2-1)^2}\Bigr)
\]
by \eqref{eq:lambda-p}. The average defining $\beta_p(h)$ may equally be taken over $a\bmod p^2$, since its summand has period $p$. Therefore, by \eqref{eq:cp-sum},
\[
(p-1)\beta_p(h)+(p^2-p)\beta_{p^2}(h)=\frac1{p^2}\sum_{a\bmod p^2}\bigl(c_p+c_{p^2}\bigr)(hf(a))=p^2\vartheta_p(h)-1.
\]
Distributing $c^2=\prod_p(1-p^{-2})^2$ over the primes, the local factor becomes
\[
\Bigl(1-\frac1{p^2}\Bigr)^2\Bigl(1+\frac{p^2\vartheta_p(h)-1}{(p^2-1)^2}\Bigr)=\frac{(p^2-1)^2+p^2\vartheta_p(h)-1}{p^4}=1-\frac2{p^2}+\frac{\vartheta_p(h)}{p^2}=w_p(h),
\]
by \eqref{eq:deltap}. No restriction on $f$ modulo $p$ has been used.
\end{proof}

\section{A partial summation lemma}

\begin{lemma}\label{lem:periodic}
Let $b$ be a periodic sequence with integer period $r\ge1$, $|b(k)|\le1$, and mean $m_b=r^{-1}\sum_{a=1}^rb(a)$. Let $1\le u\le v$ be real and let $F$ be continuously differentiable on $[u,v]$. Then
\[
\Bigl|\sum_{u<k\le v}b(k)F(k)-m_b\int_u^vF(t)\dd t\Bigr|\le4r\Bigl(|F(v)|+\int_u^v|F'(t)|\dd t\Bigr).
\]
\end{lemma}

\begin{proof}
Let $P(t)=\sum_{k\le t}b(k)-m_bt$. Writing $t=qr+s$ with $q=\lfloor t/r\rfloor$ and $0\le s<r$, the $q$ complete periods cancel, the remaining at most $r$ terms have modulus at most $1$, and $|m_bs|\le r$; so $|P(t)|\le2r$. Hence $D_u(t)=\sum_{u<k\le t}b(k)-m_b(t-u)=P(t)-P(u)$ satisfies $|D_u(t)|\le4r$ and $D_u(u)=0$. Riemann--Stieltjes integration by parts gives the identity
\[
\sum_{u<k\le v}b(k)F(k)-m_b\int_u^vF(t)\dd t=\int_u^vF(t)\dd D_u(t)=F(v)D_u(v)-\int_u^vF'(t)D_u(t)\dd t,
\]
from which the bound follows.
\end{proof}

With $F=w=f'/f$ on $[z,y]$, \cref{lem:poly}(b) gives $|w(y)|\le U/y\le U/z$ and $\int_z^y|w'|\le V/z$. Since $\int_z^yw=\log(Y/Z)$, we obtain for every $r$-periodic $b$ as in \cref{lem:periodic} and $1\le z\le y$:
\begin{equation}\label{eq:periodic-w}
\Bigl|\sum_{z<k\le y}b(k)\frac{f'(k)}{f(k)}-m_b\log\frac YZ\Bigr|\le\frac{4D_wr}z,\qquad D_w=U+V.
\end{equation}

\section{The averaging proposition}

\begin{proposition}\label{prop:averaging}
There is $C(f)>0$ such that the following holds uniformly for real $x\ge2$, $1\le z\le y$ with $Z=f(z)\le Y=f(y)\le x$, and real $2\le R\le x$. With
\[
\mathcal A=\mathcal A(x,Z,R)=Z^{-1/d}R^{1/2}(\log2R)^2+R^{-1/2}(\log2R)^2\log(2x)
\]
and
\begin{equation}\label{eq:If}
I_f(x;Z,Y)=\int_Z^Y\sum_{1\le h\le x/t}(x-ht)W_f(h)\dd t,
\end{equation}
we have
\begin{align}
\bigl|T_f(z,y)-\widetilde R_f\log(Y/Z)\bigr|&\le C(f)\,\mathcal A,\label{eq:avg1}\\
\bigl|M_f(x;z,y)-I_f(x;Z,Y)\bigr|&\le C(f)\,x^2\mathcal A.\label{eq:avg2}
\end{align}
No condition on the discriminant, content or factorisation of $f$ is imposed.
\end{proposition}

Throughout the proof, $\mathcal T(R)$ and $\mathcal N(R)$ are as in \cref{prop:second-moment}(d), so that $\mathcal T(R)\le C_0R^{-1/2}(\log2R)^2$ and $\mathcal N(R)\le C_0R^{1/2}(\log2R)^2$. If $z=y$ both assertions are trivial, so we assume $z<y$.

\begin{proof}[Proof of \eqref{eq:avg1}]
By \cref{prop:second-moment}(a) applied to $q=f(k)$,
\[
T_f(z,y)=\sum_{z<k\le y}\frac{f'(k)}{f(k)}\sum_{r\mid f(k)}\lambda(r).
\]
The part with $r>R$ is nonnegative and at most $\mathcal T(R)\sum_{z<k\le y}f'(k)/f(k)\le2U\mathcal T(R)\log(2x)$ by \eqref{eq:sum-w}. For each $r\le R$, the sequence $b(k)=\one_{r\mid f(k)}$ has period $r$ and mean $\rho(r)/r$, so \eqref{eq:periodic-w} gives
\[
\sum_{z<k\le y}\frac{f'(k)}{f(k)}\one_{r\mid f(k)}=\frac{\rho(r)}r\log\frac YZ+e_r,\qquad|e_r|\le\frac{4D_wr}z.
\]
Multiplying by $\lambda(r)$ and summing over $r\le R$ produces an error at most $4D_w\mathcal N(R)/z$. Completing $\sum_{r\le R}\lambda(r)\rho(r)/r$ to $\widetilde R_f$ costs at most $\mathcal T(R)\log(Y/Z)\le\mathcal T(R)\log(2x)$, because $\rho(r)/r\le1$ and $2\le Z\le Y\le x$. Since $Z\le\kappa_2z^d$, we have $z^{-1}\le\kappa_2^{1/d}Z^{-1/d}$, and therefore
\[
\bigl|T_f(z,y)-\widetilde R_f\log(Y/Z)\bigr|\le C_0\Bigl(4D_w\kappa_2^{1/d}Z^{-1/d}R^{1/2}(\log2R)^2+(2U+1)R^{-1/2}(\log2R)^2\log(2x)\Bigr).\qedhere
\]
\end{proof}

For the second assertion we first need a bound for the weight functions. Let $H_0=\lfloor x/Z\rfloor$. For $1\le h\le H_0$ define
\[
v_h=\min\{y,\,f^{-1}(x/h)\},\qquad F_h(t)=f'(t)\bigl(x-hf(t)\bigr)\quad(z\le t\le v_h),
\]
which is legitimate because $x/h\ge Z=f(z)$ (\cref{lem:poly}(a)). Then $z\le v_h\le(x/(\kappa_1h))^{1/d}$ and $0\le x-hf(t)\le x$ on $[z,v_h]$. Put
\[
K_2=\begin{cases}\kappa_2+\kappa_2^2/\kappa_1,&d=1,\\ \kappa_2+\kappa_2/(d-1)+\kappa_2^2/((2d-1)\kappa_1),&d\ge2.\end{cases}
\]

\begin{lemma}\label{lem:weight-var}
For $1\le h\le H_0$, $\displaystyle|F_h(v_h)|+\int_z^{v_h}|F_h'(t)|\dd t\le K_2\,x\,v_h^{d-1}$.
\end{lemma}

\begin{proof}
We have $|F_h(v_h)|\le\kappa_2xv_h^{d-1}$ and $|F_h'|\le|f''|(x-hf)+hf'^2$. For $d\ge2$,
\[
\int_z^{v_h}|F_h'|\le\kappa_2x\int_z^{v_h}t^{d-2}\dd t+h\kappa_2^2\int_z^{v_h}t^{2d-2}\dd t\le\frac{\kappa_2}{d-1}xv_h^{d-1}+\frac{h\kappa_2^2}{2d-1}v_h^{2d-1},
\]
and $hv_h^d\le x/\kappa_1$ bounds the last term by $\kappa_2^2xv_h^{d-1}/((2d-1)\kappa_1)$. For $d=1$, $f''=0$ and $\int_z^{v_h}|F_h'|\le h\kappa_2^2v_h\le\kappa_2^2x/\kappa_1$.
\end{proof}

\begin{proof}[Proof of \eqref{eq:avg2}]
\emph{Truncation in $r$.} Put
\[
\Sing_R(n)=\sum_{r\le R}\lambda(r)\frac{c_r(n)}{\varphi(r)},\qquad W_R(h)=\sum_{r\le R}\lambda(r)\beta_r(h).
\]
By \cref{prop:ramanujan-sing,lem:W-ident} and the bounds $|c_r(n)|\le\varphi(r)$ and $|\beta_r(h)|\le1$,
\[
|\Sing(n)-\Sing_R(n)|\le\mathcal T(R),\qquad|W_f(h)-W_R(h)|\le\mathcal T(R)
\]
uniformly in $n,h$. Let $M_R$ and $I_R$ be $M_f$ and $I_f$ with $\Sing$ and $W_f$ replaced by $\Sing_R$ and $W_R$. For fixed $k$ there are at most $x/f(k)$ values of $h$, each with $0\le x-hf(k)\le x$, so by \eqref{eq:sum-w}
\[
|M_f-M_R|\le\mathcal T(R)x^2\sum_{z<k\le y}\frac{f'(k)}{f(k)}\le2U\mathcal T(R)x^2\log(2x).
\]
Similarly, for each $t\in[Z,Y]$ there are at most $x/t$ values of $h$, so $|I_f-I_R|\le\mathcal T(R)x^2\int_Z^Y\dd t/t\le\mathcal T(R)x^2\log(2x)$.

\emph{The finite approximations.} All remaining sums are finite and may be rearranged. For integers $k\in(z,y]$, the condition $h\le x/f(k)$ is equivalent to $k\le v_h$, and $k>z$ forces $h<x/Z$. Hence
\[
M_R=\sum_{r\le R}\lambda(r)\sum_{h\le H_0}\sum_{z<k\le v_h}F_h(k)\frac{c_r(hf(k))}{\varphi(r)}.
\]
The sequence $k\mapsto c_r(hf(k))/\varphi(r)$ has period $r$, because $f$ has integer coefficients, and its mean is $\beta_r(h)$. The substitution $t=f(u)$ gives
\[
\int_z^{v_h}F_h(u)\dd u=\int_Z^{\min\{Y,x/h\}}(x-ht)\dd t,
\]
and therefore $I_R=\sum_{r\le R}\lambda(r)\sum_{h\le H_0}\beta_r(h)\int_z^{v_h}F_h(u)\dd u$. By \cref{lem:periodic,lem:weight-var}, the inner difference for fixed $r$ and $h$ is at most $4rK_2xv_h^{d-1}\le4rK_2\kappa_1^{-(d-1)/d}x^{2-1/d}h^{-1+1/d}$. If $v_h=z$ the difference is zero. By \cref{lem:sums}(e), $\sum_{h\le H_0}h^{-1+1/d}\le d(x/Z)^{1/d}$, so for fixed $r$ the total difference is at most $K_3rx^2Z^{-1/d}$ with $K_3=4dK_2\kappa_1^{-(d-1)/d}$. No logarithm arises here. Multiplying by $\lambda(r)$ and summing over $r\le R$ gives
\[
|M_R-I_R|\le K_3x^2Z^{-1/d}\mathcal N(R).
\]

Combining the three estimates proves
\[
|M_f-I_f|\le C_0x^2\Bigl(K_3Z^{-1/d}R^{1/2}(\log2R)^2+(2U+1)R^{-1/2}(\log2R)^2\log(2x)\Bigr).
\]
Together with the first part, \cref{prop:averaging} holds with $C(f)=C_0\max\{4D_w\kappa_2^{1/d},K_3,2U+1\}$.
\end{proof}

\begin{remark}
The error $x^2Z^{-1/d}R^{1/2}$ comes from the $k$-sampling of a weight of total variation about $xv_h^{d-1}$, summed over $h$ with the $h$-dependent endpoint $v_h$. The archived plan points out that using the endpoint $y$ for every $h$, instead of $v_h$, would lose the claimed range. This error is small only when $Z^{1/d}$, which is comparable to $z$, is large compared with $R^{1/2}(\log2R)^2$. The tail term $R^{-1/2}(\log2R)^2\log(2x)$, on the other hand, needs $R^{1/2}$ to be large compared with $(\log2R)^2\log(2x)$. The choice $R=Z^{1/d}$ in \cref{ch:assembly} balances the two powers: both terms are then $Z^{-1/(2d)}$ up to logarithmic factors.
\end{remark}
